% part: intuitionistic-logic
% chapter: introduction
% section: constructive-reasoning

\documentclass[../../../include/open-logic-chapter]{subfiles}

\begin{document}

\olfileid{int}{int}{cr}

\section{Panalaran Konstruktif}

Beda karo pangambakan logika klasik nganggo operator modal utawa
kuantor orde kapindho, logika intuisionistik kalebu ``nonklasik''
amarga mbatesi logika klasik. Logika klasik iku \emph{ora konstruktif}
ing maneka cara. Logika intuisionistik dirancang kanggo nggambarake
panalaran kang luwih ``konstruktif'' lan dadi ciri sawijining jinis
matematika konstruktif. Tuladha ing ngisor iki nggambarake sawetara
alasan dhasare.

Upamane ana wong ngaku wis nemtokake bilangan asli~$n$ kang nduweni
sipat mangkene: yen $n$ genep, hipotesis Riemann bener; yen $n$
gasal, hipotesis Riemann ora bener. Kabar apik!{} Apa hipotesis
Riemann bener utawa ora iku salah siji pitakonan gedhe kang durung
kapecah ing matematika, lan kayane wong iku wis ngowahi prakara iki
dadi prakara itungan: nemtokake apa sawijining bilangan tartamtu
genep utawa ora.

Pira nilai ajaib saka~$n$ iku? Wong iku nerangake mangkene: $n$ iku
bilangan asli kang padha karo $2$ yen hipotesis Riemann bener, lan
padha karo $3$ yen ora.

Kanthi nesu, sampeyan njaluk dhuwit bali. Saka sudut pandang klasik,
andharan mau pancen nemtokake siji nilai $n$ kang unik; nanging kang
sampeyan karepake yaiku nilai $n$ kang diwenehake kanthi
\emph{gamblang}.

Kanggo tuladha liyane kang mbokmenawa ora digawe-gawe banget,
gatekna pitakonan iki. Kita ngerti yen bilangan irasional bisa
dipangkatake nganggo bilangan rasional lan asile rasional. Contone,
$\sqrt{2}^2 = 2$. Kang durung cetha yaiku apa bilangan irasional
bisa dipangkatake nganggo bilangan \emph{irasional} lan ngasilake
bilangan rasional. Teorema iki menehi wangsulan ya:

\begin{thm}
Ana bilangan irasional $a$ lan $b$ kang ndadekake $a^b$ rasional.
\end{thm}

\begin{proof}
Gatekna $\sqrt{2}^{\sqrt{2}}$. Yen iki rasional, bukti wis rampung:
kita bisa milih $a = b = \sqrt{2}$. Yen ora, bilangan iku irasional.
Mula kita duwe
\[
(\sqrt{2}^{\sqrt{2}})^{\sqrt{2}} = \sqrt{2}^{\sqrt{2} \cdot
  \sqrt{2}} = \sqrt{2}^2 = 2,
\]
kang rasional. Dadi, ing kasus iki, pilih $a$ minangka
$\sqrt{2}^{\sqrt{2}}$ lan $b$ minangka~$\sqrt 2$.
\end{proof}

Apa iki bukti kang sah? Umume matematikawan nganggep sah. Nanging,
isih ana bab kang kurang marem: kita wis mbuktekake anane sapasang
bilangan nyata kanthi sipat tartamtu, nanging durung bisa nyebutake
\emph{pasangan endi} kang dimaksud. Asil kang padha bisa dibuktekake
kanthi nyebutake pasangan $a$, $b$ \emph{ing} bukti: pilih
$a = \sqrt{3}$ lan $b = \log_3 4$. Mula
\[
a^b = \sqrt{3}^{\log_3 4} = 3^{1/2 \cdot \log_3 4} = (3^{\log_3
  4})^{1/2} = 4^{1/2}= 2,
\]
amarga $3^{\log_3 x} = x$. Bilangan logaritma mau irasional: yen
rasional, ana pangkat bulat positif saka telu kang padha karo
pangkat bulat positif saka papat, kang mokal miturut faktorisasi prima.

Logika intuisionistik dirancang kanggo nggambarake jinis panalaran
kang ora ngidini langkah kaya ing bukti kapisan. Mbuktekake anane
$x$ kang nyukupi~$!A(x)$ ateges kudu menehi sawijining~$x$ tartamtu
lan bukti yen dheweke nyukupi $!A$, kaya ing bukti kapindho.
Mbuktekake yen $!A$ utawa $!B$ bener mbutuhake bukti kanggo salah
sijine.

Miturut cara formal, logika intuisionistik dipikolehi kanthi mbatesi
sistem !!a{derivation} kanggo logika klasik kanthi cara tartamtu.
Saka sudut matematika, iki mung sistem deduktif formal, nanging kaya
kang wis kasebut, sistem iki dirancang kanggo nggambarake sawijining
jinis panalaran matematika. Kita bisa nganggep panalaran iki sah
miturut sawijining pandangan filosofis babagan matematika (kayata
intuisionisme Brouwer); bisa uga nganggep panalaran iki luwih
``konkret'' lan luwih marem (kaya konstruktivisme Bishop); lan bisa
uga ngrembug apa andharan formal iku bener-bener nggambarake
motivasi informal. Nanging, apa wae pandangan filosofis kita,
logika intuisionistik tetep bisa ditliti minangka logika kang
disajekake kanthi formal; lan kanthi maneka alasan, akeh ahli
logika matematika kang nganggep panlitene narik kawigaten.

\end{document}
